\begin{textsample}{POS Dim 1 – 2011\_c6 – Score -2.00 – 2011\_c6/enade\_2011\_tec\_constr\_f1a5d330d7.txt}  \label{ex:f1_pos_215}
30 CAPÍTULO 3 ANÁLISE TÉCNICA DA PROVA Este capítulo tem por objetivo apresentar o desempenho dos estudantes concluintes de Tecnologia em Construção de Edifícios no ENADE/2011.
Para isso, foram calculadas as estatísticas básicas da prova em seu todo, bem como as estatísticas dos componentes relacionadas à Formação Geral, ao de Conhecimento Específico da Área e das questões discursivas isoladamente.
Nas tabelas, são apresentados o tamanho da população inscrita e de presentes, e as seguintes estatísticas das notas8 : média do desempenho na prova, erro padrão da média, desvio padrão, nota mínima, mediana e nota máxima.
As estatísticas apresentadas neste capítulo contemplam o total de estudantes concluintes da área de Tecnologia em Construção de Edifícios em 2011 do Brasil e, separadamente, por Grande Região.
Foram calculadas tendo -se em vista as seguintes agregações : ( a ) as Grandes Regiões e o país como um todo ; ( b ) a Categoria Administrativa ; e ( c ) a Organização Acadêmica.
Em relação aos gráficos de distribuição de notas, o intervalo considerado foi de 10 unidades, aberto à esquerda e fechado à direita, com exceção do primeiro intervalo, [ 0 ; 10 ], fechado em ambos os extremos.
Para os gráficos de distribuição das notas das questões discursivas, foram consideradas mais duas categorias : questão em branco e nota zero. 3.1 ESTATÍSTICAS BÁSICAS DA PROVA 3.1.1 Estatísticas Básicas Gerais A Tabela 3.1 apresenta as estatísticas básicas da prova por grande Região.
A população total de inscritos foi de 591.
Destes, 482 estiveram presentes, sendo 18,4 \% o índice de não comparecimento.
A Região de maior abstenção foi a Norte ( 28,3 \% ) e a de menor abstenção foi a Sul ( 0,0 \% ).
A média das notas da prova como um todo ( nas seções seguintes serão analisados os componentes de Formação Geral e de Conhecimento Específico ) foi 47,8, sendo que os alunos da região Nordeste obtiveram a média mais baixa ( 43,6 ) e os da região Sudeste obtiveram a média mais alta ( 51,7 ).
As demais médias foram : 44,0 na região Norte, 51,4 na região Sul e 50,1 na região Centro-Oeste.
O desvio padrão para o Brasil como um todo foi 8 Essas estatísticas e outras estão definidas no Capítulo 1. 13,1, sendo o maior desvio padrão encontrado na região Sul ( 17,6 ) e o menor na região Centro-Oeste ( 8,8 ), indicando uma menor dispersão das notas desta última região.
A região que obteve a maior nota máxima foi a Sudeste ( 80,5 ), ao passo que a região que atingiu a menor nota máxima foi a Centro-Oeste ( 66,0 ).
A mediana do Brasil como um todo foi 48,4, sendo a maior mediana obtida nas regiões Centro-Oeste ( 52,3 ) e a menor obtida nas regiões Norte e Nordeste ( 43,6 ).
A nota mínima foi zero somente na região Norte e a maior nota mínima foi obtida na região Centro-Oeste ( 33,8 ).
O comportamento das notas dos estudantes de todo o Brasil pode ser observado no Gráfico 3.1, que apresenta um histograma com a distribuição das mesmas.
Essa é uma distribuição unimodal com moda no intervalo ( 40 ; 50 ].
Apesar do coeficiente de assimetria da distribuição das notas ser negativo ( – 0,20 ) este é pequeno e podemos considerar que a distribuição é aproximadamente simétrica.
As distribuições por Grande Região também apresentam assimetria negativa, concentração pouco maior do lado direito do histograma e mais espalhada do lado esquerdo.
A assimetria negativa é mais acentuada na região Norte, onde o coeficiente encontrado foi – 1,19.
A única exceção é a região Nordeste, com coeficiente de assimetria positivo ( 0,04 ), com o comportamento oposto : concentração à esquerda e cauda um pouco mais longa à direita.
Os Gráficos 3.2, 3.3 e 3.4 apresentam informações referentes à média da nota final dos Participantes, desagregando os resultados de acordo com, respectivamente, as Grandes Regiões do país, a Categoria Administrativa e a Organização Acadêmica.
Os gráficos apresentam o valor da média das notas como uma barra e os extremos do intervalo de confiança de 95 \% como linhas verticais unidas por uma linha horizontal na forma da letra H maiúscula.
Considerando -se o gráfico de notas segundo Grande Região ( Gráfico 3.2 ), observa -se que existe diferença estatisticamente significativa, ao nível de 95 \%, entre a maior média, obtida na região Sudeste ( 51,7 ) e as obtidas nas regiões Nordeste ( 43,6 ) e Norte ( 44,0 ).
Levando -se em conta os estudantes por Categorias Administrativas ( Gráfico 3.3 ), observa -se que não há diferença estatisticamente significativa entre as médias das notas das IES Públicas ( 48,5 ) e Privadas ( 44,0 ).
Tendo como foco o Gráfico 3.4, que apresenta as notas médias das provas segundo a Organização Acadêmica, contata -se que existe diferença estatisticamente significativa ao nível de 95 \% nas médias das notas dos estudantes provenientes de Universidades e Faculdades.
Não houve alunos de Tecnologia em Construção de Edifícios provenientes de Centros Universitários.
A maior nota média foi obtida pelos estudantes de Faculdades ( 51,5 ) e a menor pelos de Universidades ( 45,8 ). 3.1.2 Estatísticas Básicas no Componente de Formação Geral A Tabela 3.2 apresenta as estatísticas básicas em relação ao componente da prova que avalia a Formação Geral dos estudantes concluintes.
Os alunos de todo Brasil obtiveram desempenho médio de 53,0.
Quanto à variabilidade, o desvio padrão das notas dos estudantes do Brasil como um todo foi 17,4.
A maior média foi obtida na região Norte ( 55,0 ), e a menor, na região Nordeste ( 51,2 ).
As demais médias foram : 54,6 na região Sudeste ; 52,9 na região Sul ; e 51,5 na região Centro-Oeste.
Já o maior desvio padrão foi obtido na região Norte ( 19,9 ), e o menor, na região Sudeste ( 16,3 ).
Os demais desvios padrões foram : 17,9 na região Nordeste ; 19,2 na região Sul ; e 17,8 na região Centro-Oeste.
A maior nota no componente de Formação Geral da prova do ENADE foi obtida por pelo menos um aluno da região Sudeste ( 92,5 ), enquanto que a menor nota máxima foi obtida na região Sul ( 77,0 ).
Nas outras regiões as notas máximas foram : 88,5 na região Norte ; 90,5 na região Nordeste ; e 78,0 na região Centro-Oeste.
A mediana do Brasil como um todo foi 55,5, sendo a menor mediana encontrada na região Nordeste ( 51,5 ) e a maior encontrada na região Norte ( 62,5 ).
A nota mínima nesta parte foi zero nas regiões Norte e Nordeste.
Nas outras regiões a nota mínina foi : 7,5 na região Sudeste ; 15,0 na região Sul ; e 7,5 na região Centro-Oeste.
O Gráfico 3.5 propicia a avaliação do desempenho dos estudantes no componente de Formação Geral a partir do histograma da distribuição das notas correspondentes.
A distribuição é unimodal, com moda em ( 60 ; 70 ], enquanto na prova como um todo a moda foi alcançada no intervalo ( 40 ; 50 ].
Nota -se, ainda, que no gráfico 3.5 as notas apresentam uma maior dispersão do que no Gráfico 3.1 ( distribuição das notas da prova ), confirmado pela comparação dos desvios padrões : 13,1 para a nota da prova como um todo e 17,4 para o componente de Formação Geral.
Para o componente de Formação Geral, o coeficiente de assimetria da distribuição das notas dos estudantes, como na prova como um todo, é negativo ( – 0,48 ).
Este coeficiente explica a concentração à direita e cauda maior à esquerda que se observa no gráfico.
Em todas as Grandes Regiões os histogramas também possuem assimetria negativa, variando entre – 0,99 ( Sul ) e – 0,34 ( Nordeste ).
Nos Gráficos 3.6, 3.7 e 3.8 são apresentadas as informações referentes ao desempenho dos Concluintes no componente de Formação Geral, em diferentes agregações : Grande Região do país, Categoria Administrativa e Organização Acadêmica.
Observa -se pelo Gráfico 3.6 que não existe diferença estatisticamente significativa entre as médias das notas no Componente de Formação Geral, segundo Grande Região do país.
Vemos que alguns dos intervalos de confiança são bem largos, em especial o da região Sul.
Este fato está relacionado com o tamanho da população envolvida nestas regiões : 15 alunos na região Sul, 26 na região Centro-Oeste e 33 na região Norte ( ver Tabela 2.3 ).
No Gráfico 3.7, que representa as notas médias no Componente de Formação Geral segundo a Categoria Administrativa do país, observa -se que existe diferença estatisticamente significativa entre as médias.
Os concluintes das IES Públicas obtiveram uma média maior ( 54,1 ) do que os das IES Privadas ( 47,0 ).
Considerando -se o tipo de Organização Acadêmica, nota -se, no Gráfico 3.8, que não há diferença estatisticamente significativa entre as médias dos alunos de Universidades ( 52,9 ) e Faculdades ( 53,4 ). 3.1.3 Estatísticas Básicas do Componente de Conhecimento Específico A Tabela 3.3 apresenta as estatísticas básicas referentes ao componente de Conhecimento Específico da área de Tecnologia em Construção de Edifícios.
A média do desempenho dos alunos do Brasil como um todo foi 46,0.
A maior média foi obtida na região Sul ( 50,9 ), e a menor, na região Norte ( 40,3 ).
As demais médias foram : 41,0 na região Nordeste ; 50,8 na região Sudeste ; e 49,6 na região Centro-Oeste.
Quanto à variabilidade das notas, o desvio padrão do Brasil como um todo foi 14,9, sendo o maior desvio padrão observado na região Sul ( 19,3 ) e o menor na região Centro-Oeste ( 10,4 ).
Os demais desvios foram : 14,6 na região Norte ; 14,1 na região Nordeste ; e 14,2 da região Sudeste.
A mediana das notas dos estudantes de todo o Brasil foi 46,4.
A maior mediana ocorreu na região Sul ( 50,8 ) e a menor na região Nordeste ( 39,9 ).
As demais medianas foram : 41,9 na região Norte ; 50,3 na região Sudeste ; e 47,9 na região Centro-Oeste.
A nota máxima do Brasil como um todo foi 83,0, obtida por pelo menos um aluno da região Sul.
A menor nota máxima ocorreu na região Centro-Oeste, 66,6.
As demais notas máximas foram : 69,0 nas regiões Norte ; 75,0 na região Nordeste ; e 81,0 na região Sudeste.
A nota mínima foi zero apenas na região Norte.
Nas demais regiões as notas mínimas foram : 8,7 na Nordeste ; 11,6 na Sudeste ; 16,0 na Sul ; e 28,0 na Centro-Oeste.
Assim como os Gráficos 3.1 e 3.5, o Gráfico 3.9, apresentado a seguir, proporciona uma avaliação do desempenho de concluintes em relação ao componente de Conhecimento Específico com um histograma da distribuição das notas correspondentes.
Dentre as três distribuições apresentadas, esta é a mais simétrica ( coeficiente de assimetria de – 0,09 ).
Esta também é uma distribuição unimodal, e o grupo modal é o ( 40,50 ].
Os coeficientes de assimetria das distribuições das notas do componente de Conhecimento Específico das regiões Norte, Sudeste e Centro-Oeste são negativos, – 0,51 ; – 0,34 e – 0,11, respectivamente.
Já nas regiões Nordeste ( 0,18 ) e Sul ( 0,07 ), os coeficientes de assimetria são positivos e bem pequenos, evidenciando que a distribuição das notas dos estudantes dessas regiões são aproximadamente simétricas.
Os Gráficos 3.10, 3.11 e 3.12 apresentam uma comparação dos resultados em relação à Grande Região do país, à Categoria Administrativa e à Organização Acadêmica, agora levando em conta o desempenho de alunos no componente de Conhecimento Específico da prova.
Pelo Gráfico 3.10, observa -se que uma das diferenças estatisticamente significativa é a existente entre as regiões Sudeste ( 50,8 ) e a Norte ( 40,3 ), mesmo esta última apresentando um intervalo de confiança largo.
Quanto à Categoria Administrativa ( gráfico 3.11 ), observa -se um comportamento semelhante ao da prova como um todo, ou seja, não existe diferença estatisticamente significativa entre as médias das IES Públicas ( 46,6 ) e Privadas ( 43,0 ) no componente de Conhecimento Específico.
Quanto ao Gráfico 3.12, observa -se que existe diferença estatisticamente significativa ao nível de 95 \% entre as notas no componente de Conhecimento Específico das Universidades ( 43,4 ) e Faculdades ( 50,8 ). 3.2 ANÁLISE DAS QUESTÕES OBJETIVAS 3.2.1 Componente de Formação Geral A Tabela 3.4 apresenta as estatísticas básicas relativas às oito questões objetivas do componente da prova que abrange a Formação Geral dos estudantes.
A média do Brasil foi 53,5.
A menor média foi encontrada na região Sul ( 46,7 ) e a maior na região Norte ( 60,2 ).
As demais médias foram : 51,9 na região Nordeste ; 54,6 na região Sudeste ; e 51,4 na região Centro-Oeste.
O desvio padrão do Brasil foi 19,0, sendo o maior desvio padrão encontrado na região Norte ( 21,3 ) e o menor na região Centro-Oeste ( 17,8 ).
Os demais desvios foram : 18,5 na região Nordeste ; 19,0 na região Sudeste ; e 19,7 na região Sul.
A mediana das notas do Brasil como um todo foi 50,0, a mesma das regiões Nordeste, Sul e Centro-Oeste.
Nas regiões Norte e Sudeste a mediana foi maior, 62,5.
A maior nota máxima 100,0 foi atingida nas regiões Nordeste e Sudeste.
A menor nota máxima, 75,0, ocorreu nas regiões Sul e Centro-Oeste.
As notas mínimas nas regiões Norte, Nordeste e Sudeste foi 0,0 ; e nas regiões Sul e Centro-Oeste, 12,5.
A Tabela 3.5 apresenta os índices de facilidade e discriminação ( ponto bisserial ) das questões objetivas do componente de Formação Geral.
Quanto ao índice de facilidade, foram usadas as seguintes cores para diferenciar o nível de dificuldade da questão : • Azul para as questões classificadas com índice muito fácil ( índice \&gt ; =0,86 ), verde para as questões classificadas com índice fácil ( 0,61 a 0,85 ), amarelo para as questões classificadas com médio ( 0,41 a 0,60 ), vermelho para as questões classificadas com difícil ( 0,16 a 0,40 ) e roxo para as questões classificadas com muito difícil ( \&lt ; =0,15 ).
Já quanto ao índice de discriminação, foram usadas as seguintes cores para qualificar a questão : • As questões classificadas com índice fraco receberam a cor vermelho ( índice \&lt ; =0,19 ), as classificadas com médio receberam a cor amarelo ( 0,20 a 0,29 ), as classificadas com bom receberam a cor verde ( 0,30 a 0,39 ) e as classificadas com muito bom ( \&gt ; =0,40 ) receberam a cor azul.
As questões objetivas do componente de Formação Geral, segundo o índice de facilidade, foram assim avaliadas : das oito questões, nenhuma teve o índice de facilidade classificado como muito fácil.
Quatro questões foram tidas como fáceis, por terem índice de acertos situado na faixa entre 0,61 e 0,85 ( de 61,0 \% a 85,0 \% de acertos ).
Duas questões foram consideradas de dificuldade média, situando -se no intervalo entre 0,41 e 0,60 do índice de facilidade, ou seja, houve entre 41,0 \% e 60,0 \% de acertos.
Outras duas questões foram classificadas na categoria difícil, situando -se no intervalo entre 0,16 e 0,40.
Por fim, nenhuma das questões apresentou menos de 15 \% de acertos, razão pela qual seriam classificadas como muito difíceis.
Como já comentado, para análise das questões objetivas relativas à Formação Geral segundo o poder de discriminação, utilizou -se, o índice de discriminação ponto bisserial.
Nesta análise as questões foram assim avaliadas : exatamente metade das oito questões apresentaram índices acima ou igual a 0,40 e, assim, foram classificadas com índice muito bom para esse grupo de estudantes.
A outra metade teve bom índice de discriminação, com valor entre 0,30 e 0,39.
Nenhuma questão foi classificada nem com índice médio, nem com nível fraco de discriminação.
O índice de facilidade variou de 0,17 a 0,80, e o de discriminação, de 0,32 a 0,59.
As quatro questões com índices de discriminação muito bom, as de números 1, 2, 3 e 5, figuraram entre os diversos níveis de dificuldade desse conjunto : três classificadas na categoria fácil ( questões 1, 3 e 5 ) do índice de facilidade e outra na categoria médio ( questão 2 ).
Em particular, a questão 3 foi a que apresentou o maior poder discriminatório, com índice 0,59, e foi também uma das mais fáceis, com uma proporção de 0,71 acertos.
O máximo de acertos foi alcançado pela questão 6 com um índice de facilidade de 0,80.
A questão de número 4 apresentou um índice de facilidade de 0,52, ou seja, um pouco mais da metade dos estudantes conseguiu resolvê -la, dentro do universo de participantes.
Seu índice de discriminação foi bom ( 0,39 ).
Já a questão 8 obteve índice de discriminação bom, 0,33, e seu índice de facilidade foi difícil ( 0,17 ).
O Gráfico 3.13, para exemplificar, analisa o comportamento da questão de número 3 de Formação Geral.
Trata -se da segunda questão mais fácil e a que obteve o maior índice de discriminação dessa parte da prova.
Neste gráfico, cada uma das cinco curvas representa o percentual de respostas em determinada alternativa da questão, em função do número de acertos dos estudantes nesta parte da prova ( Formação Geral/Múltipla Escolha ), antes de possíveis eliminações pelo critério do ponto bisserial.
Em princípio, a soma das escolhas possíveis deveria ser igual a 100 \%.
Não é este o caso, notadamente para o número zero de acertos, pois não aparecem no gráfico as questões deixadas em branco ou com múltiplas respostas.
Como foram oito as questões, os valores variam de 0 a 8 acertos.
A curva em vermelho corresponde à alternativa E, a correta para esta questão.
Assim, observa -se que entre os estudantes com menor número de acertos, nessa parte do exame, a situação mais frequente foi a escolha de uma das alternativas incorretas : a alternativa A ( em azul ), B ( em verde ) ou D ( em roxo ).
Na medida em que o número de acertos aumenta, indicando desempenho melhor nesta parte da prova, aumenta concomitantemente a proporção de estudantes que selecionaram a alternativa correta E, atingindo 100 \% para os estudantes com 7 acertos.
Essa análise permite verificar como a questão discriminou os grupos de desempenho, justificando o alto índice obtido na questão.
Os gráficos relativos às demais questões de Formação Geral constam do Anexo I. 3.2.2 Componente de Conhecimento Específico A Tabela 3.6 apresenta as estatísticas básicas em relação às questões objetivas do componente de Conhecimento Específico da prova, por Grande Região.
A média do Brasil deste componente foi de 50,2.
A menor média foi observada na região Norte ( 45,2 ) e a maior na região Centro-Oeste ( 54,7 ).
O desvio padrão de todo o Brasil foi 16,0, sendo o menor desvio padrão encontrado na região Centro-Oeste ( 12,2 ) e o maior na região Sul ( 20,8 ).
A mediana de todo o Brasil foi 50,0.
A menor mediana ( 45,5 ) ocorreu nas regiões Norte e Nordeste e a maior mediana foi 54,5, das regiões Sudeste e Sul.
A nota máxima das questões objetivas do componente de Conhecimento Específico foi 90,9, obtida por pelo menos um aluno das regiões Sudeste e Sul.
As notas máximas nas demais regiões foram : 72,7 na região Norte ; 81,8 na região Nordeste ; e 77,3 na região Centro-Oeste.
A Tabela 3.7 apresenta os índices de facilidade e discriminação ( ponto bisserial ) das questões objetivas do componente de Conhecimento Específico.
Para facilitar a diferenciação das questões usaram -se as mesmas cores da Tabela 3.5 para as diferentes classificações dos índices de facilidade e de discriminação.
Dentre as questões objetivas da parte da prova relativa ao componente de Conhecimento Específico, nenhuma foi anulada pela Comissão.
Desse modo, a classificação quanto ao índice de facilidade foi estabelecida com base em todas as 27 questões.
A partir dos índices obtidos, pode -se concluir que as questões objetivas da prova estavam mais concentradas entre as questões difíceis, com 11 questões nesta categoria ( a categoria modal ).
As outras 16 questões estavam distribuídas com valores decrescentes a partir desta : nove classificadas como de dificuldade média e sete como fácil.
Não houve questão classificada como muito difícil, nem tampouco muito fácil.
Já quanto aos índices de discriminação das questões objetivas do componente de Conhecimento Específico da prova, tem -se como resultado a seguinte classificação : oito das 27 questões foram consideradas como boas, enquanto 4 delas tiveram índice de discriminação muito bom.
Dentre as demais, dez delas foram classificadas como médias e cinco como fracas, sendo quinze, por conseguinte, a quantidade de questões nos dois patamares mais baixos de discriminação.
Constata -se, assim, que a prova – no que se refere ao componente de Conhecimento Específico – possuía capacidade fraca de discriminar entre aqueles que dominam ou não o conteúdo.
Dentre as quatro questões que alcançaram os maiores índices de discriminação, as de números 23 e 26 foram classificadas como fáceis e as outras duas ( 10 e 29 ) com grau de dificuldade médio.
A Tabela 3.7 apresenta os índices de facilidade e discriminação ( ponto bisserial ) das questões objetivas do componente de Conhecimento Específico.
Para facilitar a diferenciação das questões usaram -se as mesmas cores da Tabela 3.5 para as diferentes classificações dos índices de facilidade e de discriminação.
Dentre as questões objetivas da parte da prova relativa ao componente de Conhecimento Específico, nenhuma foi anulada pela Comissão.
Desse modo, a classificação quanto ao índice de facilidade foi estabelecida com base em todas as 27 questões.
A partir dos índices obtidos, pode -se concluir que as questões objetivas da prova estavam mais concentradas entre as questões difíceis, com 11 questões nesta categoria ( a categoria modal ).
As outras 16 questões estavam distribuídas com valores decrescentes a partir desta : nove classificadas como de dificuldade média e sete como fácil.
Não houve questão classificada como muito difícil, nem tampouco muito fácil.
Já quanto aos índices de discriminação das questões objetivas do componente de Conhecimento Específico da prova, tem -se como resultado a seguinte classificação : oito das 27 questões foram consideradas como boas, enquanto 4 delas tiveram índice de discriminação muito bom.
Dentre as demais, dez delas foram classificadas como médias e cinco como fracas, sendo quinze, por conseguinte, a quantidade de questões nos dois patamares mais baixos de discriminação.
Constata -se, assim, que a prova – no que se refere ao componente de Conhecimento Específico – possuía capacidade fraca de discriminar entre aqueles que dominam ou não o conteúdo.
Dentre as quatro questões que alcançaram os maiores índices de discriminação, as de números 23 e 26 foram classificadas como fáceis e as outras duas ( 10 e 29 ) com grau de dificuldade médio.
A questão de número 28 foi a mais difícil dentre as 27 questões específicas, classificada como difícil pelo índice de facilidade com apenas 16,0 \% de acertos.
Essa questão apresentou poder discriminatório fraco, 0,15.
Pelo critério do ponto bisserial esta questão foi eliminada do cômputo da nota final.
Além dessa, outras quatro questões, as de números 17, 25, 30 e 32 apresentaram também baixo poder discriminatório, classificadas como fracas pelo índice de discriminação e foram igualmente eliminadas do cômputo da nota final.
Estas questões foram todas classificadas como difíceis com respeito ao índice de facilidade.
A título de exemplo das análises do comportamento das questões objetivas, o Gráfico 3.14 analisa a questão 23 do componente de Conhecimento Específico.
Esta foi a uma das questões mais fáceis da prova, apresentando índice de facilidade 0,62, ou seja, 62,0 \% dos estudantes assinalaram acertadamente a opção D, correspondente ao gabarito.
Seu índice de discriminação foi igual a 0,45, classificado como muito bom.
Neste gráfico, cada uma das cinco curvas representa o percentual de respostas em determinada alternativa da questão 23 em função do número de acertos dos estudantes nesta parte da prova, antes de possíveis eliminações de questões pelo critério do ponto bisserial.
A alternativa correta D, representada no gráfico pela curva em roxo, foi escolhida em maiores proporções pelos alunos com desempenho melhor nesta parte da prova.
Já as alternativas incorretas, também denominadas distratores, foram selecionadas principalmente por aqueles com notas mais baixas.
Em princípio a soma dos estudantes que escolheram cada uma das alternativas deveria ser igual a 100 \%, mas alguns dos estudantes com nota zero deixaram esta questão em branco ou marcaram mais de uma alternativa, comportamento considerado inválido.
A proporção de alunos que selecionou a resposta correta D aumenta gradativamente, chegando a atingir 100 \% para mais de 20 acertos, enquanto a proporção dos que escolheram alternativas incorretas, principalmente a alternativa C – curva em preto, decai a partir de quatro acertos, como função do número de acertos nesta parte da prova.
Os gráficos relativos às demais questões do componente de Conhecimento Específico constam do Anexo I. 3.3 ANÁLISE DAS QUESTÕES DISCURSIVAS 3.3.1 Componente de Formação Geral As análises dos resultados de desempenho dos estudantes de Tecnologia em Construção de Edifícios nas duas questões discursivas relativas à Formação Geral encontram -se na Tabela 3.8 e no Gráfico 3.15.
Na tabela 3.8 observa -se que as notas médias dos estudantes de todo o Brasil nesse conjunto de questões não difere da obtida nas questões objetivas.
A média nas questões objetivas foi 53,5 ( com erro padrão 0,9 ) e nas questões discursivas foi 52,4 ( com erro padrão 1,2 ).
No entanto, há aumento do desvio padrão de 19,0, nas questões objetivas do componente de Formação Geral dos alunos de todo o Brasil, para 27,2 nas questões discursivas do mesmo componente.
Esse comportamento – valor próximo das médias e aumento no desvio padrão – pode ser observado nas regiões Nordeste, Sudeste e Centro- O Gráfico 3.15 representa a distribuição das notas nas questões discursivas no componente de Formação Geral.
A moda desta distribuição ocorre no intervalo ( 40 ; 50 ].
Observa -se que neste conjunto de questões – discursivas de Formação geral – mais alunos obtiveram notas acima do intervalo modal, o que se confirma pelo valor da mediana ( 55,0 ) e do coeficiente de assimetria ( – 0,40 ).
Os coeficientes de assimetria de todas as regiões são negativos, sendo o da região Sul o de maior valor absoluto ( – 1,32 ) e o da região Norte o de menor valor absoluto ( – 0,10 ).
Oeste.
Diferentemente, na região Norte o desempenho nas questões discursivas ( 47,3 ) foi menor do que nas objetivas ( 60,2 ).
Já na região Sul, a média obtida nas questões discursivas de Formação Geral ( 62,2 ) foi maior do que a obtida nas objetivas deste mesmo componente ( 46,7 ).
A mediana de todo o Brasil, neste componente, foi 55,0, sendo a maior mediana encontrada na região Sul ( 70,0 ) e a menor nas regiões Norte e Nordeste ( 50,0 ).
A nota máxima ( 100,0 ) não foi obtida por alunos das regiões Sul e Centro-Oeste, onde a máxima foi 90,0.
A nota mínima ( 0,0 ) foi a mesma em todas as regiões.
O Gráfico 3.15 representa a distribuição das notas nas questões discursivas no componente de Formação Geral.
A moda desta distribuição ocorre no intervalo ( 40 ; 50 ].
Observa -se que neste conjunto de questões – discursivas de Formação geral – mais alunos obtiveram notas acima do intervalo modal, o que se confirma pelo valor da mediana ( 55,0 ) e do coeficiente de assimetria ( – 0,40 ).
Os coeficientes de assimetria de todas as regiões são negativos, sendo o da região Sul o de maior valor absoluto ( – 1,32 ) e o da região Norte o de menor valor absoluto ( – 0,10 ).
Na sequência, os resultados verificados para cada uma das questões discursivas de Formação Geral serão apresentados, estabelecendo -se relações com os conteúdos abordados em cada uma delas.
Os comentários da Banca de docentes corretores a respeito do observado na correção das respostas dos estudantes, suas impressões e conclusões serão apresentados junto à análise de cada questão.
Cumpre esclarecer que, tendo em vista que as questões discursivas de Formação Geral são padronizadas, ou seja, constam de todas as provas, os comentários da Banca são os mesmos para todas as carreiras acadêmicas, sendo direcionados a todos os estudantes que participaram do ENADE/2011.
A seguir, serão analisados os desempenhos da Área de Tecnologia em Construção de Edifícios nas duas questões discursivas de Formação Geral do ENADE/2011, comparando os resultados obtidos com comentários para cada questão. 3.3.1.1 Análise da Questão Discursiva 1 do Componente de Formação Geral Os dados de Tecnologia em Construção de Edifícios, obtidos a partir das respostas à questão 1, encontram -se na Tabela 3.9 e no Gráfico 3.16.
Nessa questão – de melhor desempenho dentre as duas de Formação Geral – os alunos de todo Brasil tiveram média, 56,8.
A maior média para a questão 1 foi obtida na região Sul ( 71,7 ), e a menor, na região Norte ( 51,7 ).
Quanto à variabilidade das notas, o desvio padrão de todo o Brasil foi 33,9.
O menor desvio padrão foi obtido na região Sudeste ( 31,2 ) e o maior desvio padrão foi obtido na região Norte ( 37,4 ).
As medianas das notas dos estudantes de todo o Brasil foi 60,0, a mesma das regiões Nordeste e Sudeste.
A menor mediana ocorreu na região Norte ( 40,0 ) e a maior na região Sul ( 90,0 ).
A nota máxima 100,0 e a nota mínima 0,0 foram iguais em as regiões do Brasil.
O Gráfico 3.16 mostra a distribuição das notas na questão discursiva 1 do componente de Formação Geral.
Observa -se que a maior frequência corresponde aos alunos que obtiveram notas maiores que 90,0.
Também se observa concentração de notas no intervalo ( 20 ; 30 ].
O coeficiente de assimetria é – 0,21 indicando leve assimetria negativa.
Dentre as Grandes Regiões apenas na Norte o coeficiente de assimetria é positivo, mas como é bem pequeno ( 0,04 ), pode -se considerar que a distribuição das notas desta região é virtualmente simétrica. 3.3.1.2 Comentários sobre a correção das respostas à Questão Discursiva 1 De maneira geral, a aplicação da chave de correção da questão 1 de Formação Geral não apresentou qualquer dificuldade digna de menção.
Para isso contribuíram, indubitavelmente, os ajustes feitos após a correção da amostra e a reunião entre todos os membros da Banca de docentes corretores.
As poucas dúvidas, todas pontuais, apresentadas pelos corretores, foram acompanhadas e respondidas pela coordenação e subcoordenação da correção das questões de Formação Geral, por meio da ferramenta de Gerenciamento de Dúvidas do Sistema de Correção On-line.
Não houve registro de qualquer ocorrência que pusesse em xeque o padrão de resposta ou a efetividade e a adequação da chave de correção.
Explica -se : trata -se de questão com comando claro, direto e objetivo ( solicitava -se, basicamente, três vantagens justificadas de cursos a distância ), cujas respostas foram corrigidas por meio da aplicação de um chave de correção testada e aprovada previamente.
Havia absoluta clareza quanto aos critérios de avaliação da correspondência entre as respostas dos estudantes e as possibilidades de vantagens de cursos a distância admitidas como corretas no padrão de resposta oficial, além de gradações explícitas ( e fáceis de aplicar ) dos diferentes níveis de pontuação previstos.
Felizmente, portanto, não há reparo a registrar em relação à facilidade de aplicação do padrão de resposta e da chave de correção, e nem em relação à atribuição dos diferentes níveis de pontuação previstos.
Todas as dificuldades que poderiam ter obstado a correta aplicação do padrão de resposta oficial e da respectiva chave de correção foram evitadas por meio dos ajustes feitos após a correção da amostra e fartamente debatidos com toda a Banca.
Digno de nota é que quantidade tão significativa de profissionais envolvidos na mesma tarefa – tanto para a questão 1 quanto para a 2 – tenha apresentado tão poucas dificuldades na execução da correção, em termos proporcionais.
Em suma, a correção da questão 1 da prova de Formação Geral do ENADE 2011 foi exemplarmente bem planejada, servindo -se de padrão de resposta muito bem adequado à questão proposta.
Quanto ao tema desta questão, em particular, a Banca verificou que uma parcela significativa de estudantes evocou experiências bastante concretas e próximas de sua realidade.
Houve várias respostas que indicavam uma vivência pessoal de ensino superior na modalidade Educação a Distância ( EaD ), evidenciando o tom de depoentes nos textos apresentados.
Os estudantes, em sua grande maioria, utilizaram parte considerável do espaço de 15 linhas disponíveis para a resposta – e outra parcela menos significativa dos que não o fizeram demonstraram notável capacidade de atender ao comando da questão de maneira objetiva, curta e, via de regra, correta.
Registre -se, ainda, que foram relativamente poucos os casos de respostas que tenham passado ao largo do tema em pauta na questão.
As capacidades de leitura, de compreensão do comando proposto e de expressão escrita que os estudantes avaliados na edição 2011 do ENADE foram satisfatórias.
Não obstante, seria leviano perder de vista que a qualidade dos textos redigidos em resposta às questões discursivas do Exame ainda está muito aquém do que se espera de concluintes de cursos de ensino superior de todas as regiões do país.
Quanto ao conteúdo das respostas, a Banca constatou boa capacidade, por parte da maioria dos estudantes, de compreensão do tema e do comando da questão.
Foram relativamente poucos os casos de respostas que deixaram de enumerar vantagens da modalidade EaD, e proporcionalmente escassos os estudantes que citaram vantagens não previstas no padrão de resposta.
Foi frequente, a tentativa direta de atender ao comando da questão.
Os erros mais comuns, em relação ao padrão de respostas e à grade de correção, foram fruto do desdobramento em vários “ itens ” daquilo que, de acordo com o padrão de respostas oficial, representava uma única vantagem.
Destacaram -se, neste caso, as respostas que apontavam a flexibilidade de horário e/ou local como duas vantagens distintas daquela modalidade de ensino.
Quanto aos diferentes níveis de pontuação previstos, a maior causa de baixas pontuações foi a ausência de justificativas, e mesmo de argumentação, para uma ou mais das vantagens enumeradas.
Isso demonstra não apenas a objetividade das respostas, coerente com a objetividade do comando da questão ( “ enumere três vantagens de um curso a distância ” ), mas também certa dificuldade de formulação plena de um texto, ou ao menos de parágrafos, em formato dissertativo – mesmo diante de uma média de 5 linhas disponíveis para cada vantagem a enumerar.
Os acertos mais comuns, ou seja, os “ itens ” do padrão de respostas mais frequentemente mencionados foram : ( 1 ) a flexibilidade de horário e/ou local ; ( 2 ) a capilaridade do ensino a distância ; ( 3 ) a democratização do acesso à educação de qualidade ; e ( 4 ) os custos menores que os de cursos presenciais.
Dentre as vantagens previstas no padrão de respostas que foram menos citadas, destacam -se a inclusão de pessoas com comprometimento motor, a qualificação de professores e a troca de experiências entre os participantes.
Foram muito comuns, no entanto, as menções à supostamente maior facilidade de acesso a professores ou/e tutores em cursos superiores a distância. 3.3.1.3 Análise da Questão Discursiva 2 do Componente de Formação Geral A Tabela 3.10 mostra que o desempenho dos estudantes de Tecnologia em Construção de Edifícios na questão 2.
A média 48,0 foi inferior à obtida na questão de número 1 ( 56,8 ).
A região Sul foi aquela onde a média, nessa questão, foi maior ( 52,7 ), e a de menor média foi a da região Norte ( 42,9 ).
Quanto à variabilidade das notas, o desvio padrão de todo o Brasil foi 31,1, inferior ao obtido na questão de número 1 ( 33,9 ).
O maior O maior desvio nessa questão foi obtido na região Norte ( 36,3 ), enquanto o menor foi obtido na região Sudeste ( 29,7 ).
A mediana comprova a diferença de desempenho entre essas duas regiões : na primeira questão foi 60,0 e na segunda questão foi 50,0.
Em todas as regiões a mediana foi 50,0 para esta questão, com exceção da região Sul, onde esta estatística foi 60,0.
A nota máxima ( 100,0 ) foi obtida em todas as regiões do Brasil, bem como a nota mínima 0,0.
O Gráfico 3.17 mostra a distribuição das notas na questão discursiva 2 do componente de Formação Geral.
Observa -se que o percentual de alunos que deixou esta questão em branco, além de superior ao da questão 1, é a moda desta distribuição.
Desconsiderando -se os alunos que não resolveram a questão, o intervalo de maior frequência é aquele ( 50 ; 60 ].
Como na questão 1, o coeficiente de assimetria é negativo, – 0,37.
Em todas as regiões a assimetria é negativa, mais acentuada na região Sul ( – 0,78 ) e mais leve na região Norte ( – 0,06 ). 3.3.1.4 Comentários sobre a correção das respostas à Questão Discursiva 2 Diferentemente da questão 1, a questão 2, cujo tema aborda políticas públicas para a erradicação do analfabetismo e para o aumento da empregabilidade, permitiu mais amplitude na elaboração das respostas dos estudantes, uma vez que solicita a apresentação de uma proposta.
À falta de objetividade técnica na elaboração da questão – sobretudo no seu comando – correspondeu uma miríade de possibilidades interpretativas.
Registre -se que não raras foram as respostas que se utilizaram de lugares comuns e exortações religiosas/humanitárias/cívicas, contudo, na grande maioria dos casos, essas exortações foram usadas como um encerramento do texto e não comprometeram a resposta.
Muito frequente também foi a confusão entre política pública e política partidária, bem como alguns poucos, confundiam programas educacionais com programas educativos nas TVs.
Assim, destacaram -se como propostas/programas mais recorrentes :  prosseguimento das políticas já existentes, mas com o aumento dos investimentos, normalmente sugerindo a ampliação da bolsa família, da Educação de Jovens e Adultos ( EJA ) e da EaD ;  valorização do magistério, construção de escolas e melhoria das atuais ;  parceria do governo com empresas para manutenção de salas de aula visando a alfabetização de seus funcionários, oferecimento de estágios e redução da carga horária em troca de isenção de impostos ;  parcerias com igrejas e ONGs para criação de espaços de alfabetização ;  escolas itinerantes e alfabetizadores em domicílio, principalmente para pessoas com dificuldade de locomoção, como os idosos e deficientes físicos ;  erradicação do trabalho infantil ;  vinculação da bolsa família não apenas à frequência, mas também e, principalmente, ao resultado obtido pelo aluno na escola ;  revisão das políticas atuais, sendo a mais frequente a extinção da aprovação automática ;  críticas consistentes em relação a modelos didáticos considerados inadequados e desestimulantes para a educação não só de adultos, mas de pessoas de todas as idades.
Destacam -se, ainda, outras sugestões apresentadas :  Creches nas escolas onde os pais estão sendo alfabetizados.  Diminuição de duas horas na jornada de trabalho em empresas para funcionários não alfabetizados, para que possam frequentar a escola.  Campanhas educativas vinculadas aos meios de comunicação.  Aumento do número de escolas noturnas.  Formação específica para professores alfabetizadores.
Quanto à relação entre o analfabetismo e a empregabilidade, deve -se sublinhar que nem todos os estudantes estabeleceram claramente o vínculo entre essas duas situações sociais.
Alguns falaram separadamente de uma e de outra.
Mas a maioria fez referência à necessidade de estudo para “ conseguir um bom emprego com um bom salário ”.
Alguns, em menor número, estabeleceram de forma bastante interessante a questão histórica para a situação do Nordeste ; e, também, a relação entre a pessoa analfabeta/com pouca instrução e a desempregada/empregada em serviços mais pesados e pior remunerados, que não teriam condições de educar seus filhos que, por sua vez, também não teriam melhores oportunidades no mercado de trabalho, identificando a formação de um círculo vicioso e a necessidade de sua interrupção.
Muitos afirmaram que o analfabetismo não é o único responsável pelo desemprego e sim a má distribuição de renda.
Grande também foi o número de alunos que criticaram o resultado da pesquisa, afirmando que a mesma é enganosa, uma vez que considera alfabetizada a pessoa que “ desenha seu nome ”, sendo comum associarem a atual política de alfabetização com ganhos eleitorais.
Foi comum a resposta incompleta, em que o estudante fez apenas a análise das desigualdades/crítica do quadro apresentado, ou só apresentou proposta.
Alguns se limitaram a responsabilizar o governo referindo -se de forma bastante genérica à questão da “ educação ” e “ profissional ”.
Entre os equívocos que mais se repetiram, destacou -se a simples análise dos dados apresentados na tabela que consta do enunciado.
Alguns textos, inclusive, estavam corretos, sem que, no entanto, fosse respondida a pergunta.
Da mesma forma, quando os estudantes partiram para a segunda parte da questão, fizeram referência a vários programas já existentes ou simplesmente disseram que é muito importante que existam projetos para a educação.
Outro equívoco recorrente foi a análise da educação no Brasil como um todo.
Essa análise, apesar de correta, não se referia especificamente à questão do analfabetismo, o que, por conseguinte, levava a sugestões que não eram direcionadas à erradicação do analfabetismo, e sim à melhora da educação no Brasil.
Sendo assim, foram apresentadas propostas como o aumento do número de faculdades, o reforço de alunos do Ensino Médio, a criação de escolas técnicas, etc.
Por outro lado, houve um grande número de redações bem escritas e precisas, no que se refere ao que foi exigido pela questão.
Foram análises equilibradas e sensatas, correlacionando a problemática do analfabetismo com o desemprego e a apresentação de sugestões bastante consistentes. 3.3.2 Componente de Conhecimento Específico Na parte da prova relativa às questões discursivas no componente de Conhecimento Específico ( Tabela 3.11 ), observa -se que a média foi mais baixa do que para as questões discursivas do componente de Formação Geral.
Enquanto no componente de Formação Geral a média para estudantes de Tecnologia em Construção de Edifícios de todo o Brasil foi 52,4, na parte de Conhecimento Específico a média foi 22,6.
A maior média deste componente foi obtida pelos estudantes da região Sul ( 31,9 ), e a menor, pelos da região Norte ( 12,4 ).
Quanto à variabilidade das notas, o desvio padrão de todo o Brasil foi 16,1.
O maior desvio padrão foi encontrado na região Sul ( 16,8 ), e o menor, na região Norte ( 11,8 ).
A maior nota máxima foi obtida na região Sudeste ( 70,0 ) enquanto a menor nota máxima foi encontrada na região Centro-Oeste ( 43,3 ).
Além disso, a nota mínima ( 0,0 ) foi obtida por alunos de todas as regiões do Brasil, com exceção da região Sul ( 3,3 ).
A mediana do Brasil como um todo foi 21,7, também encontrada na região Centro-Oeste.
As demais medianas foram : 10,0 na região Norte ; 15,0 na região Nordeste ; e 31,7 nas regiões Sudeste e Sul.
O Gráfico 3.18 representa a distribuição das notas nas questões discursivas no componente de Conhecimento Específico.
Essa distribuição é unimodal com moda no intervalo de [ 0,10 ].
Nota -se que a distribuição é decrescente a partir do intervalo ( 20 ; 30 ] até o intervalo ( 60 ; 70 ].
Não foram registradas notas maiores que 70,0.
O coeficiente de assimetria é 0,38. componente foi 22,6.
A maior média deste componente foi obtida pelos estudantes da região Sul ( 31,9 ), e a menor, pelos da região Norte ( 12,4 ).
Quanto à variabilidade das notas, o desvio padrão de todo o Brasil foi 16,1.
O maior desvio padrão foi encontrado na região Sul ( 16,8 ), e o menor, na região Norte ( 11,8 ).
A maior nota máxima foi obtida na região Sudeste ( 70,0 ) enquanto a menor nota máxima foi encontrada na região Centro-Oeste ( 43,3 ).
Além disso, a nota mínima ( 0,0 ) foi obtida por alunos de todas as regiões do Brasil, com exceção da região Sul ( 3,3 ).
A mediana do Brasil como um todo foi 21,7, também encontrada na região Centro-Oeste.
As demais medianas foram : 10,0 na região Norte ; 15,0 na região Nordeste ; e 31,7 nas regiões Sudeste e Sul.
O Gráfico 3.18 representa a distribuição das notas nas questões discursivas no componente de Conhecimento Específico.
Essa distribuição é unimodal com moda no intervalo de [ 0,10 ].
Nota -se que a distribuição é decrescente a partir do intervalo ( 20 ; 30 ] até o intervalo ( 60 ; 70 ].
Não foram registradas notas maiores que 70,0.
O coeficiente de assimetria é 0,38. 3.3.2.1 Análise da Questão Discursiva 3 do Componente de Conhecimento Específico Na questão 3, cujos resultados aferidos encontram -se descritos na Tabela 3.12, a média dos estudantes de todo o Brasil foi 28,8.
A menor média nessa questão foi obtida pelos alunos da região Norte ( 13,6 ), enquanto a maior média foi obtida na região Sudeste ( 45,6 ).
Quanto à variabilidade das notas, o desvio padrão de todo o Brasil foi 29,9.
O maior desvio padrão foi obtido na região Sudeste ( 32,9 ), enquanto o menor foi obtido na região Norte ( 17,6 ).
A nota máxima, 100,0 pontos, foi alcançada por pelo menos um aluno da região Sudeste.
A região onde ocorreu a menor nota máxima foi a Norte ( 60,0 ).
A mediana do Brasil como um todo foi 25,0.
Nas regiões Norte e Nordeste a mediana foi 0,0, indicando que mais de 50 \% dos alunos destas regiões deixaram a questão em branco ou obtiveram nota zero.
Nas demais regiões as medianas foram : 40,0 na Sudeste ; 30,0 na sul ; e 10,0 na Centro-Oeste.
A nota mínima ( 0,0 ) foi obtida em todas as regiões do Brasil, sem exceção.
O Gráfico 3.19 mostra a distribuição das notas na questão discursiva 3, do componente de Conhecimento Específico.
Essa distribuição é \textbf{bimodal} com moda na categoria questões “ em branco ” e no intervalo ( 20 ; 30 ].
Os coeficientes de assimetria da distribuição de notas de todos os alunos de Tecnologia em Construção de Edifícios do Brasil é positivo ( 0,79 ).
O mesmo ocorre para a distribuição de todas as regiões, sendo o da região Sudeste o menor ( 0,06 ), e o da região Nordeste o maior ( 1,03 ). 3.3.2.2 Comentários sobre a correção das respostas à Questão Discursiva 3 A questão é pertinente ao dia a dia das construções, envolvendo elementos básicos do planejamento e controle de obras.
O tema deve ser de conhecimento de quem pretende acompanhar/coordenar atividades de construção, utilizando redes de planejamento/procedência e cronogramas de obras.
O padrão de resposta apresenta as informações/elementos que devem constar no cronograma em questão e os pontos essenciais da resposta, em função dos quais foram atribuídos valores para cada nível de conhecimento apresentado pelos concluintes.
Grande parte dos estudantes respondeu apenas dois ou três itens elementares, ficando no nível I de conhecimento apresentado no padrão de resposta.
Outros citaram mais elementos, porém muito poucos conseguiram listar o conjunto elencado no padrão de resposta.
Nem sempre as respostas foram apresentadas em forma de tópicos, porém, a partir dos textos, foi possível avaliar o nível de conhecimento do formando.
Elementos não pertinentes a esse tipo de cronograma foram desconsiderados.
Observou -se que, dentre os estudantes que responderam, a maior parte se limitou às informações mais elementares, como descrição das atividades, tempo de duração ou início e fim de cada atividade.
Outros citaram o início e o fim da obra, mas não de cada atividade.
Outro grupo apresentou respostas medianas, que envolviam o conceito de caminho crítico e folga das atividades.
Poucos apresentaram a ideia global de folgas livre e total, o caminho crítico e a possibilidade de reestudo de prazos das atividades em relação ao prazo para o término da obra.
Observou -se, ainda, que as respostas muitas vezes englobavam elementos pertencentes a outros tipos de controle, notadamente o que se refere a controle financeiro e quantidade de materiais e mão de obra em cada atividade.
Pelas respostas apresentadas, a maioria dos concluintes não sabe exatamente os elementos que compõem especificamente um cronograma físico PERT-CPM, pois os elementos mais apresentados pertencem ao cronograma solicitado ( físico PERT-CPM ), como também a cronogramas físicos de modo geral.
Outro aspecto relevante é a confusão com elementos que pertencem a outros tipos de controle como, por exemplo, a cronogramas financeiros. 3.3.2.3 Análise da Questão Discursiva 4 do Componente de Conhecimento Específico A Tabela 3.13 contém as informações relativas à questão 4 do conjunto de questões do componente de Conhecimento Específico.
O desempenho dos estudantes de todo o Brasil nesta questão foi inferior ao desempenho na questão de número 3.
A média geral do Brasil foi 20,6, sendo a menor média registrada na região Norte ( 10,8 ) e a maior na região Sul ( 42,0 ).
A maior nota máxima foi 80,0, atingida por pelo menos um aluno da região Sudeste.
A menor nota máxima ( 55,0 ) ocorreu na região Norte.
A mediana do Brasil como um todo ( 15,0 ) foi a mesma da região Nordeste.
Na região Norte a mediana foi 0,0 e a maior mediana ocorreu na região Sul ( 45,0 ).
A nota mínima ( 0,0 ) foi a mesma em todas as regiões do Brasil, sem exceção.
O Gráfico 3.20, representa a distribuição de notas na questão discursiva 4, do componente de Conhecimento Específico.
O intervalo ( 20 ; 30 ] foi o de maior frequência, mas observa -se outra concentração na categoria “ em branco ”.
Comparando -se com os dados da questão discursiva 3, houve uma maior proporção de estudantes que deixaram essa questão em branco.
O coeficiente de assimetria deste histograma ( Gráfico 3.20 ) é positivo ( 0,58 ), como também são os das regiões Norte ( 1,24 ), Nordeste ( 0,87 ) e Sudeste ( 0,39 ).
Já nas regiões Sul e Centro-Oeste a assimetria é negativa, respectivamente, – 0,65 e – 0,13. 3.3.2.4 Comentários sobre a correção das respostas à Questão Discursiva 4 A questão aborda elementos que devem ser observados na especificação de materiais para compra, bem como os cuidados a serem observados quando da sua aplicação/instalação.
Esquadrias de alumínio estão presentes em boa parte das construções de edifícios no Brasil, sendo, portanto, relevante o conhecimento de suas características por parte dos responsáveis pela compra de materiais e/ou por sua instalação durante a obra.
A questão foi dividida em duas partes, uma relacionada à especificação e outra à instalação, abordando conceitos relacionados a gerenciamento de execução de serviços e técnicas construtivas.
Em relação ao item “ a ” da questão, a maior parte dos estudantes listou apenas um ou dois elementos dos apresentados no padrão de resposta esperado.
O item “ b ” teve um número menor de respostas e boa parte dos alunos apresentou cuidados a serem observados na colocação dos contramarcos na alvenaria ( como prumo, nível, esquadro ) e não os cuidados na colocação das esquadrias nos contramarcos, como apresentado no padrão de resposta.
Observou -se que, no item “ a ”, a maioria dos que responderam parcialmente à questão citou as dimensões como elemento a ser especificado.
Poucos citaram o tipo de perfil a ser utilizado.
Quanto ao item “ b ”, a grande maioria dos que responderam citou elementos com o prumo e alinhamento referentes à colocação do contramarco, o que não está contemplado no padrão de resposta.
Provavelmente não deve ter ficado claro para o estudante, neste caso, que o pretendido na questão eram os cuidados na colocação da esquadria em si, partindo da ideia que o contramarco já estaria colocado.
Em função das respostas apresentadas, é possivel deduzir que os concluintes entenderam o que foi soliciatdo no item “ a ”, o que pode não ter ocorrido no item “ b ”.
Todos os elementos elencados no padrão de resposta foram citados por um ou outro estudante, apesar de a maioria ter se limitado a poucos elementos.
Ressalta -se, ainda, a grande quantidade de questões em branco. 3.3.2.5 Análise da Questão Discursiva 5 do Componente de Conhecimento Específico A Tabela 3.14 contém as informações relativas à questão 5 do conjunto do componente de Conhecimento Específico.
O desempenho dos estudantes nessa questão foi inferior ao das questões 3 e 4.
A nota média dos estudantes de todo o Brasil foi 18,4.
A maior média foi registrada na região Sul ( 28,0 ), enquanto a menor média foi registrada na região Norte ( 12,9 ).
Quanto à variabilidade das notas, o desvio padrão dos alunos do Brasil, como um todo, foi 14,7.
Enquanto o maior desvio foi encontrado na região Sul ( 16,7 ), o menor foi encontrado na região Nordeste ( 13,6 ).
A nota máxima ( 75,0 ) foi alcançada nas região Sul, enquanto a menor nota máxima foi registrada nas regiões Norte e Centro-Oeste ( 50,0 para ambas ).
A mediana para todo o Brasil foi 20,0, e a nota mínima foi zero para quase todas as regiões do Brasil, com exceção da região Sul, onde a nota mínima foi 5,0.
O Gráfico 3.21 representa a distribuição das notas da questão discursiva 5 no componente de Conhecimento Específico.
O número de estudantes que deixaram a questão 5 em branco é próximo ao que ocorreu com a questão 4, e a moda encontra -se no intervalo inferior, ( 10 ; 20 ].
Os coeficientes de assimetria são positivos, variando entre 0,27 ( Sudeste ) e 1,58 ( Sul ), sendo 0,42 o do Brasil como um todo. 3.3.2.6 Comentários sobre a correção das respostas à Questão Discursiva 5 A questão aborda um tema atual, diretamente ligado à construção civil, principalmente pelo alto consumo de materiais e pelo elevado índice de desperdício observado nas construções de modo geral.
A busca de novos processos e materiais que diminuam a agressão ambiental deve ser um dos objetivos de quem atua no ramo da construção de edifícios.
A resposta esperada é um texto dissertativo considerando aspectos elencados nos itens “ a ”, “ b ” e “ c ” da questão.
Dentre os estudantes que responderam à questão, grande parte se limitou a comentários genéricos sobre o tema, repetindo as ideias apresentadas na introdução do item.
Em outras provas foi possível encontrar respostas que citassem, ainda que parcialmente, os elementos elencados no padrão de resposta esperado : processos/materiais que devem ser considerados na busca da sustentabilidade na construção civil, as causas do desperdício, do alto consumo de materiais e grande volume de resíduos observados na construção civil e a necessidade de novas tecnologias para a racionalização e conservação de materiais/água/energia na construção civil.
Observou -se que no item “ a ”, a maioria destacou que a construção civil é uma das maiores geradoras de resíduos e consumidoras de materiais, e que hoje em dia é crescente a preocupação com a sustentabilidade, mas não apontou caminhos para a melhoria dos índices.
No item “ b ”, os comentários foram mais superficiais, limitando -se, por vezes a repetir as palavras do enunciado (... deve -se evitar o desperdício, consumindo de forma adequada e evitando a geração de entulho ).
Pouquíssimos apontaram causas para os índices de desperdícios e resíduos.
No item “ c ” ( benefícios da racionalização ), muitos destacaram benefícios para a construtora em termos de economia ( redução dos gastos ) e de publicidade.
Outros destacaram que o maior benefício é para a natureza, pois recursos estariam sendo poupados.
Pelas respostas apresentadas, pode -se inferir que os estudantes não entenderam que, para responder à questão, era necessário citar processos ou materiais que devem ser considerados na busca da sustentabilidade, as causas do desperdício, do alto consumo de materiais e grande volume de resíduos observados na construção civil, além de citar a busca de tecnologias para a conservação de materiais, água e energia. 3.3.3 Considerações Finais sobre as questões discursivas específicas Para a banca de docentes corretores, o baixo índice de respostas com alta pontuação nas questões discursivas específicas pode ter ocorrido pela falta de conhecimento dos estudantes em relação ao tema proposto ou pela má interpretação das questões.
Outro aspecto que deve ser observado é a forma como o questionamento foi apresentado e se o padrão de resposta está de acordo com ela, pois os estudantes podem não ter alcançado os objetivos pretendidos com a questão, em função de não terem entendido exatamente o que se pretendia com o solicitado.
Observou -se ainda a dificuldade, que alguns participantes apresentaram, em desenvolver textos ou apresentar ideias.
É fundamental, portanto, que esse tipo de avaliação discursiva continue acontecendo para que sirva de estímulo ao desenvolvimento desse tipo de habilidade, tão importante para profissionais de qualquer área de atuação.
Cabe ao Tecnólogo em Construção de Edifícios não somente acompanhar a execução dos projetos, mas também gerenciar a execução dos serviços e instalações, bem como gerenciar a produção e realizar o controle de qualidade dos materiais e serviços e elaborar cronogramas e orçamentos.
Para tanto, é necessário que o profissional seja capaz de se comunicar eficientemente nas formas escrita, oral e gráfica.

% matched lemmas: bimodal
\end{textsample}
